% Generated by Nexus Scholar Harness
\documentclass{article}
\usepackage[utf8]{inputenc}
\usepackage{amsmath,amssymb}
\usepackage{cite}
\usepackage{hyperref}

\begin{document}

\section{Literature Review: Reliability and Generalization in UAV Crop-Weed Mapping - Protocol v0.3 Amendment Workspace}

\subsection{Objectives \& Research Questions}
\begin{itemize}
\item \textbf{RQ1: How do UAV crop-weed mapping studies define and evaluate reliability under spectral degradation, sensor differences, and temporal shift?}
\item \textbf{RQ2: Which datasets, sensors, dates, fields, and split units support the reported generalization claims?}
\item \textbf{RQ3: How are uncertainty, calibration, missing evidence, and failed transfer reported?}
\item \textbf{RQ4: Which evidence gaps remain after separating dataset diversity, model mechanisms, and demonstrated external validation?}
\end{itemize}

\subsection{Scope and provenance of this synthesis}

This is a \textbf{pilot evidence synthesis over 9 of the 644 included studies}, not the review
of the full corpus. The nine full texts were retrieved as Open Access PDFs on
2026-10-31, extracted to Markdown, and are listed in
extracted/\_pilot\_manifest.json (nine EXTRACTED\_VALID, one UNRESOLVED; 425,839
body characters in total). Every quoted fragment below is a verbatim substring of
one of those extracts, and the locator names the nearest extracted section heading.

Three limits must travel with this text:

\begin{enumerate}
\item \textbf{No charting rows exist yet.} The v0.3.1 charting instrument is frozen but no
   study has been charted under it, so this narrative is not a substitute for the
   protocol matrix and does not claim to be one.
\item \textbf{No dual coding, no human review of this text.} The HumanConfirmation@1 in
   literature/screening/ covers the screening decisions. It does not review,
   approve, or entail this synthesis.
\item \textbf{Extraction is block-based.} Section headings are heuristically detected, so
   locators are approximate pointers into the PDF, not page-accurate citations.
\end{enumerate}

\subsection{Evidence base}

\begin{itemize}
\item \textbf{SCI-000006} (Khan et al., 2021, journal article). Semi-supervised GAN for crop/weed
  classification on UAV RGB imagery at early growth, pea and strawberry crops. Split:
  70\% train, 15\% validation, 15\% test. Headline: 90\% average accuracy with 80\% of the
  training data unlabeled.
\item \textbf{SCI-000044} (Mekhalfa \& Yacef, 2021, preprint). Colour and texture features with an
  SVM on a public soybean UAV dataset. Split: random train/test division. Headline:
  above 96\% accuracy for combined colour and texture features.
\item \textbf{SCI-000041} (Galymzhankyzy \& Martinson, 2025, preprint). Lightweight transformer-CNN
  hybrid using RGB, NIR and red-edge bands for segmentation on the WeedsGalore maize
  benchmark. Split: 109 spatial training tiles, 23 validation, 24 test, no overlap.
  Headline: 78.88\% mIoU, 15.8 points above RGB-only.
\item \textbf{SCI-000151} (Hashemi-Beni et al., 2022, journal article). U-Net, SegNet, FCN and
  DeepLabV3+ with transfer learning for crop/weed segmentation on two open benchmark
  UAS datasets (field robot and UAV). Split: not stated in the extract. Headline:
  DeepLabV3+ 84.3\%, best of the compared architectures.
\item \textbf{SCI-000291} (Kebede et al., 2025, journal article). Fine-tuned CNN detection and
  classification of teff weeds on 1,308 drone images across growth stages and weather.
  Split: train\_test\_split, 80\% train / 20\% test. Headline: fine-tuned MobileNetV2
  96.40\% accuracy.
\item \textbf{SCI-000372} (Li \& Zou, 2026, journal article in press). UNet-ResNet segmentation
  with a NAFNet deblurring front-end on UAV sorghum imagery. Split: not stated in the
  extract. Headline: ResNet-34 mean Dice 0.9198; deeper encoders did not improve it.
\item \textbf{SCI-000607} (Sapkota et al., 2022, journal article). Mask R-CNN trained on synthetic
  images for weed detection and biomass estimation in cotton on UAV imagery. Split: not
  stated in the extract. Headline: real-plant synthetic mAPm 0.60 against 0.80 on real
  images; GAN-derived synthetic data weaker still.
\item \textbf{SCI-000644} (Zhang et al., 2021, journal article). Fastai CNN multi-label weed
  mapping in sod production, more than 10,000 images from six surveys. Split: not stated
  in the extract. Headline: precision 0.57–0.90 and recall 0.52–0.99 across five
  classes.
\item \textbf{SCI-000736} (Prabhakar et al., 2026, preprint). Cross-field and cross-crop detection
  transfer between a new cotton dataset and a soybean dataset collected under a similar
  protocol. Split: explicit source/target separation, 25-shot target budget. Headline:
  few-shot fine-tuning with 25 target examples outperformed unsupervised DAOD.
\end{itemize}

One further selected study (SCI-000227, arXiv 2109.02716) is recorded as UNRESOLVED: the
downloader built https://arxiv.org/pdf/arxiv.2109.02716 and received HTTP 404. It is
reported here rather than silently dropped.

\subsection{RQ1 — How reliability is defined and evaluated}

\textbf{Degradation is usually acknowledged, rarely isolated.} SCI-000372 is the only
study in this pilot that builds the degradation into the model rather than the
discussion: "Motion blur disproportionately affects fine-detail segmentation
accuracy by smearing edges and textures" (Section 2.5), and the paper inserts a
NAFNet restoration stage before segmentation. Its encoder comparison is a useful
negative result about capacity: "ResNet-34 achieves the highest mean Dice Score of
0.9198 and that accuracy does not increase with depth" (Results). Depth is not
reliability.

\textit{}Multispectral gains are reported against an RGB-only comparator, which is the
right shape of evidence.\textit{} SCI-000041 states that the model "achieves a
segmentation accuracy (mean IoU) of 78.88\%, outperforming RGB-only models by 15.8
percentage points" (Abstract). That is a matched-modality comparison in the sense
the protocol asks for, though it is one benchmark (WeedsGalore) and one split.

\textbf{Most remaining studies evaluate a single, favourable condition.} SCI-000291
reports "fine-tuned MobileNetV2 achieved the highest accuracy (96.40\%)" (Abstract)
after augmentation and optimizer sweeps, but the evaluation is a random image-level
80/20 split of imagery collected from the same farms. SCI-000044 reaches above
96\% accuracy with colour and texture features on a public soybean dataset, again
under a random division. SCI-000006 reports "an average accuracy of 90\% when 80\%
of training data was unlabeled" (Abstract) for a semi-supervised GAN, which is a
labelling-efficiency claim rather than a reliability-under-shift claim. These are
legitimate in-distribution results; none of them tests sensor substitution,
illumination drift, or a held-out acquisition date.

\textbf{Where the pilot does test shift, the result is a caution.} SCI-000736 states
that "Models trained under these conditions tend to generalize poorly when deployed
in new fields, seasons, or to different crops" (Abstract) and finds that "few-shot
fine-tuning with as few as 25 labeled target examples outperforms unsupervised DAOD
in our cross-crop comparison" (Abstract). The recovery came from a small amount of
target supervision and a domain-adjacent source, not from a more elaborate
adaptation algorithm.

\subsection{RQ2 — What supports the generalization claims}

The split unit is the decisive detail, and the pilot shows how unevenly it is
reported. Of the nine studies:

\begin{itemize}
\item \textbf{One explicitly spatial split}: SCI-000041 reports that "The dataset is
  spatially split into 109 training (70\%), 23 validation (15\%), and 24 test (15\%)
  tiles, ensuring no spatial overlap between splits" (Abstract). No-overlap tiles
  are stronger than a random tile split, but they remain one field and one season.
\item \textbf{Two explicitly random splits}: SCI-000291 uses "the train\_test\_split
  function, with 80\% of the data allocated for training and 20\% for testing"
  (Methodology), and SCI-000044 states that features "were divided randomly into
  two sets: the training set, and testing set" (Abstract). Both support
  in-distribution accuracy claims only.
\item \textbf{One explicit cross-domain design}: SCI-000736 separates source and target
  crops and fields, which is why its degradation finding is the most informative
  result in the pilot.
\item \textbf{The other five do not state the split unit in the extracted text}
  (SCI-000006, SCI-000151, SCI-000372, SCI-000607, SCI-000644). This is
  a reporting gap in the extract, not proof that the paper omits it; it is exactly
  the field that charting must resolve against the PDF.
\end{itemize}

Dataset diversity in the pilot is real but shallow in the generalization sense:
one maize benchmark, one soybean dataset, one cotton dataset, one sorghum set, one
teff set, one sod-farm survey series, one pea/strawberry set, and two mixed
benchmarks. Only SCI-000736 evaluates a model across two of them. SCI-000607
changes the \textit{training} distribution (synthetic images) but evaluates on real images
from the same crop, which is a data-efficiency result rather than an external
validation.

\subsection{RQ3 — Uncertainty, calibration, and failed transfer}

\textit{}No study in the nine reports a predictive-calibration metric for its weed
predictions.\textit{} There is no ECE, Brier score, reliability diagram, or predictive
uncertainty estimate in the extracted texts. The word "calibration" appears twice,
and neither occurrence is predictive calibration:

\begin{itemize}
\item SCI-000291 refers to sensor calibration before acquisition, noting the
  "importance of sensor calibration prior to data acquisition" (Related works).
\item SCI-000372 describes a training protocol in which deblurring and segmentation
  are "jointly calibrated" (Methodology), meaning joint optimisation, not
  probability calibration.
\end{itemize}

This is the terminology trap the charting codebook warns about: hardware or sensor
calibration and training co-optimisation are not uncertainty calibration. On this
pilot, the honest charting value for uncertainty\_calibration would be
not\_reported for all nine, not none and not calibrated.

\textit{}Failure is reported, but usually as a performance number rather than an explicit
failure analysis.\textit{} SCI-000607 gives the clearest negative result: synthetic
training data "did not perform as effectively as real plant instance-based synthetic
images" for GAN-derived instances, and the real-plant synthetic condition reached
mAPm 0.60 against 0.80 on real images. SCI-000736 documents cross-field
degradation and a recovery path. The remaining seven report only their best
condition and do not quantify a failure mode, which means \textit{no failure reported} must
not be read as \textit{no failure observed}.

\subsection{RQ4 — Evidence gaps after the pilot}

\begin{enumerate}
\item \textbf{Split-unit reporting is the largest gap.} Two of nine explicitly random, one
   explicitly spatial and non-overlapping, one cross-domain, and five that do not
   state the split unit in the extract. Any generalization claim that is not tied
   to a held-out field, flight, or season is currently unsupported, regardless of
   reported accuracy.
\item \textbf{Calibration and uncertainty are absent.} Zero of nine report a calibration
   metric. The protocol's uncertainty\_calibration dimension is therefore likely to
   be not\_reported for most of the corpus, and the review should say so rather than
   inferring confidence from softmax outputs.
\item \textbf{Matched-modality evidence is rare.} SCI-000041 compares multispectral fusion
   against RGB-only; the other eight do not isolate the sensor contribution, so the
   common claim that multispectral data improves weed mapping is not tested here.
\item \textbf{Synthetic-data gains are conditional.} SCI-000607 shows a large gap between
   real and synthetic training performance, which argues against treating synthetic
   augmentation as a substitute for target-domain data.
\item \textbf{The pilot is small.} Nine of 644 included studies, selected for OA
   availability, not at random. It demonstrates that the retrieval → extraction →
   synthesis path produces grounded text; it does not estimate corpus-level
   proportions, and no number in this section should be read as a PRISMA-style
   frequency over the review.
\end{enumerate}

\subsection{Method note}

Quotes were copied from the extracted Markdown and checked as verbatim substrings
(case- and whitespace-normalised) by the pilot driver's follow-up check. Study IDs
are the workspace SCI-* identities from literature/included.json. "Not stated in
the extract" means the extracted text does not contain the information; it is a
statement about this text, not about the PDF, and charting must re-check the PDF
before any dimension is recorded as not\_reported.


% No .bib file was present at export time; bibliography omitted.

\end{document}
